% III. ПРОЕКТИРАНЕ НА СИСТЕМАТА
% Тук се събира по-голямата част от покритието на учебната програма: класове,
% операции, атрибути, отговорности, CRC карти, речник на системата, връзки,
% зависимост и наследяване, ограничения, етикетирани стойности и стереотипи,
% класификатори, видимост и обхват, шаблонни класове, диаграми на случаите на
% употреба и на дейностите.
\section{Проектиране и програмна реализация}
\label{sec:design}

\subsection{Модули и посока на зависимостите}

Системата е пет модула с еднопосочни зависимости. \term{Четецът} обхожда единицата за превод и
произвежда вътрешен модел. \term{Моделът} е независим от езика на входа и от формата на изхода и е
единственото място със структура. \term{Излъчвателите} го превръщат в текст на избраната нотация,
\term{пораждащият модул} в релационна схема и в заготовки, а \term{проверяващият модул} съпоставя два
модела. Зависимостите сочат само към модела, така че нито един излъчвател не знае за съществуването
на четеца \figref{fig:packages}.

\begin{figure}[H]
\centering
\begin{tikzpicture}[node distance=1.0cm and 1.2cm]
  \node[block] (model) {\textbf{модел}\\ класификатори, членове,\\ връзки};
  \node[accent, left=2.4cm of model] (reader) {\textbf{четец}\\ AST в JSON\\ към модел};
  \node[accent, above=1.1cm of model] (diagramreader) {\textbf{четец на диаграми}\\ текст към модел};
  \node[accent, right=2.4cm of model] (emit) {\textbf{излъчватели}\\ PlantUML, Mermaid};
  \node[accent, below left=1.2cm and 0.2cm of model] (gen) {\textbf{пораждащ модул}\\ схема и заготовки};
  \node[accent, below right=1.2cm and 0.2cm of model] (check) {\textbf{проверяващ модул}\\ разлики};
  \node[note, left=1.0cm of reader, text width=2.1cm] (src) {единица\\ за превод};
  \node[note, right=1.0cm of emit, text width=2.1cm] (out) {текст на\\ диаграма};

  \draw[line] (src) -- (reader);
  \draw[line] (reader) -- (model);
  \draw[line] (diagramreader) -- (model);
  \draw[line] (model) -- (emit);
  \draw[line] (emit) -- (out);
  \draw[line] (model) -- (gen);
  \draw[line] (model) -- (check);
\end{tikzpicture}
\caption{Модули и посока на зависимостите. Всяка стрелка сочи натам, накъдето тече информация, и нито един модул не зависи от друг освен от модела}
\label{fig:packages}
\end{figure}

Правилото се проверява механично, а не се пази по споразумение: заглавните файлове на излъчвателите,
на пораждащия и на проверяващия модул включват \code{model.hpp} и не включват \code{reader.hpp}.

\subsection{Речник, класификатори и връзки}

Речникът е получен по метода на CRC картите~\cite{beck1989crc}: за всяка предполагаема отговорност е
записан носителят ѝ и сътрудниците му, след което картите без собствено поведение са отпаднали или са
се слели със съседите си. Класификаторът носи квалифицирано име, вид, шаблонни параметри, атрибути,
операции, стереотипи и етикетирани стойности; атрибутът име, тип, видимост, обхват и кратност;
операцията и признаци за виртуалност, чистота и константност. Статичните членове се маркират като
членове с обхват на класификатора. Видовете класификатори са пет: клас, структура, интерфейс, изброен
тип и шаблон. Интерфейсът не се чете от изходния текст, защото C++ няма такава ключова дума, а се
\term{извежда}: класификатор без атрибути, чиито операции са само чисти виртуални, е интерфейс по
определението на UML, и изводът се прави върху модела, а не върху синтаксиса.

Връзките се извеждат от типовете на членовете, а не се обявяват отделно. Обобщението следва пряко
от списъка на базовите класове, заедно с достъпа на всяка база и с признака за виртуално
наследяване. Зависимостта идва от употребата на тип в подпис на операция, когато няма по-силна
връзка, а реализацията свързва специализация с общия шаблон. Разграничението между асоциация,
агрегация и композиция се прави по притежанието и е изнесено в таблица, за да бъде ревизуемо
поотделно \tabref{tab:ownership}.

\begin{table}[H]
\caption{Съответствие между изписването на члена и вида на връзката}
\label{tab:ownership}
\centering
\fontsize{10}{11}\selectfont
\setstretch{1.0}
\begin{tabular}{@{}L{4.6cm}L{2.6cm}L{3.0cm}L{1.8cm}@{}}
\toprule
\textbf{Изписване на члена} & \textbf{Притежание} & \textbf{Връзка} & \textbf{Кратност} \\
\midrule
\code{T value}                    & изключително & композиция  & 1 \\
\code{std::unique\_ptr<T>}        & изключително & композиция  & 0..1 \\
\code{std::optional<T>}           & изключително & композиция  & 0..1 \\
\code{std::shared\_ptr<T>}        & споделено    & агрегация   & 0..1 \\
\code{std::weak\_ptr<T>}          & споделено    & агрегация   & 0..1 \\
\code{T*}                         & няма         & асоциация   & 0..1 \\
\code{T\&}                        & няма         & асоциация   & 1 \\
контейнер от някое от горните     & като елемента & като елемента & * \\
тип само в подпис на операция     & няма         & зависимост  & \\
\bottomrule
\end{tabular}
\end{table}

Прочитът следва UML буквално. Композиция означава, че частта загива заедно с цялото, затова член по
стойност или в указател с изключително притежание е композиция; \code{shared\_ptr} задържа частта
жива и след цялото, което е агрегация; обикновен указател или отправка не твърдят нищо за времето на
живот, което оставя проста асоциация. Член, чийто тип не е класификатор от модела, остава обикновен
атрибут и не поражда връзка. Шаблонните класове се записват като параметризирани класификатори, а
специализациите като отделни класификатори със връзка на реализация към общия шаблон. Разгръщането
на всяка инстанциация в отделен клас би затрупало диаграмата с еднотипни възли и би заличило самата
параметризация, която е носителят на смисъла; по същата причина \term{инстанциациите, породени от
употребата}, не влизат в модела, тъй като са свойство на конкретния превод, а не на проектното
решение.

\subsection{Съответствие с релационна схема и поведенчески модели}

Клас със стереотипа за постоянно съхранение поражда таблица, всеки негов атрибут със скаларен тип
поражда колона, а етикетираните стойности задават име на таблицата, име и тип на колоната и участие в
първичния ключ. Клас без обозначен първичен ключ получава сурогатен ключ \code{id}, и това се съобщава
изрично. Асоциацията с кратност едно поражда външен ключ в притежаващата страна, а с кратност много
свързваща таблица; изборът е съзнателен, защото диаграмата показва само единия край, тоест от нея не
се вижда дали връзката е едно към много или много към много, а свързващата таблица е вярна и в двата
случая.

За наследяването са реализирани и трите утвърдени стратегии~\cite{ambler2003agiledb}: една таблица за
цялата йерархия с колона разграничител, таблица на всеки клас с ключ, който е и външен ключ към базата,
и таблица на всеки конкретен клас с копирани надолу колони. По подразбиране се прилага \term{една
таблица за йерархията}, защото само при нея извличането на един обект не изисква съединяване; изборът
не е зазидан, а се прави с етикетирана стойност върху базовия клас.

Нито диаграмата на случаите на употреба \figref{fig:usecase}, нито тази на дейностите
\figref{fig:activity} може да бъде извлечена от изходен текст: случаят на употреба е твърдение за
намерение, а дейността за поток на управление, който структурата на кода не носи. Затова двете се
описват на ръка и се превеждат към същите две нотации като диаграмата на класове.

\begin{figure}[htbp]
\centering
\begin{adjustbox}{max width=\linewidth, max totalheight=0.36\textheight}
\begin{tikzpicture}[node distance=0.7cm and 2.0cm]
  \tikzset{bpuc/.style={draw, ellipse, align=center, font=\normalsize, minimum width=5.4cm,
                      minimum height=1.05cm, fill=blue!5, inner sep=2pt},
           bpact/.style={draw, align=center, font=\normalsize, minimum width=2.1cm,
                       minimum height=0.8cm, fill=green!8, rounded corners}}
  \node[bpuc] (uc1) {UC1: извличане на модел\\ от изходен текст};
  \node[bpuc, below=of uc1] (uc2) {UC2: излъчване на\\ диаграма на класове};
  \node[bpuc, below=of uc2] (uc3) {UC3: пораждане на\\ релационна схема};
  \node[bpuc, below=of uc3] (uc4) {UC4: пораждане на заготовки\\ по схема};
  \node[bpuc, below=of uc4] (uc5) {UC5: прочитане на диаграма\\ обратно в модел};
  \node[bpuc, below=of uc5] (uc6) {UC6: отчитане на разминаване\\ между диаграма и код};
  \node[bpuc, below=of uc6] (uc7) {UC7: избор на стратегия\\ за наследяването};

  \node[bpact, left=of uc2] (dev) {Разработчик};
  \node[bpact, left=of uc5] (arch) {Проектант};
  \node[bpact, below=1.4cm of arch] (ci) {Конвейер за\\ интеграция};

  \node[draw, dashed, rounded corners, fit=(uc1)(uc7), inner sep=9pt,
        label={[font=\normalsize\bfseries]above:Blueprint}] (sys) {};

  \foreach \a/\b in {dev/uc1, dev/uc2, dev/uc3, dev/uc4, arch/uc2, arch/uc5, arch/uc6, ci/uc6}
      \draw (\a) -- (\b);

  % Отношенията между случаите се извеждат вдясно от рамката на системата,
  % за да не минават през самите елипси.
  \draw[line, dashed] (uc2.east) -- ++(0.7,0)
        node[right, font=\small, xshift=1pt] {$\ll$include$\gg$} |- (uc1.east);
  \draw[line, dashed] (uc3.east) -- ++(2.6,0)
        node[above, font=\small, yshift=2pt] {$\ll$include$\gg$} |- (uc1.east);
  \draw[line, dashed] (uc6.east) -- ++(0.7,0)
        node[right, font=\small, xshift=1pt] {$\ll$include$\gg$} |- (uc5.east);
  \draw[line, dashed] (uc7.east) -- ++(4.5,0)
        node[above, font=\small, yshift=2pt] {$\ll$extend$\gg$} |- (uc3.east);
\end{tikzpicture}
\end{adjustbox}
\caption{Случаи на употреба. Изходният текст на модела е в \code{examples/blueprint.usecase} и се превежда към двете нотации с \code{blueprint usecase}}
\label{fig:usecase}
\end{figure}

\begin{figure}[htbp]
\centering
\begin{adjustbox}{max width=\linewidth, max totalheight=0.36\textheight}
\begin{tikzpicture}[node distance=0.75cm and 1.1cm]
  \tikzset{bpstep/.style={draw, rounded corners, align=center, font=\normalsize,
                        minimum width=5.2cm, minimum height=0.85cm, fill=blue!5},
           bpdec/.style={draw, diamond, aspect=2.6, align=center, font=\small,
                       fill=yellow!12, inner sep=1pt},
           bpdot/.style={draw, circle, fill=black, minimum size=5pt, inner sep=0pt}}
  \node[bpdot] (s) {};
  \node[bpstep, below=of s] (a1) {Изпълнение на фронтенда върху\\ единицата за превод};
  \node[bpstep, below=of a1] (a2) {Прочитане на синтактичното дърво\\ в модела};
  \node[bpdec, below=of a2] (d1) {има ли класификатор\\ със стереотип за съхранение};
  \node[bpstep, below left=1.0cm and 0.4cm of d1] (a3) {Излъчване на диаграмата\\ в двете нотации};
  \node[bpstep, below right=1.0cm and 0.4cm of d1] (a4) {Извеждане на\\ релационната схема};
  \node[bpstep, below=of a3] (a7) {Прочитане на диаграмата\\ обратно в модел};
  \node[bpstep, below=of a4] (a5) {Пораждане на заготовки\\ по схемата};
  \node[bpstep, below=of a5] (a6) {Извличане на модел от\\ породените заготовки};
  \node[bpstep, below=1.6cm of a7] (a8) {Съпоставяне на двата модела};
  \node[bpdec, below=of a8] (d2) {съвпадат ли};
  \node[bpstep, below left=0.9cm and 0.1cm of d2] (a10) {Отчитане на затворен цикъл};
  \node[bpstep, below right=0.9cm and 0.1cm of d2] (a9) {Отчитане на разминаване\\ и прекъсване на построяването};
  \node[bpdot, below=1.1cm of d2, yshift=-1.0cm, minimum size=6pt, label={[bpdot, minimum size=11pt, fill=none]center:}] (e) {};

  \draw[line] (s) -- (a1);
  \draw[line] (a1) -- (a2);
  \draw[line] (a2) -- (d1);
  \draw[line] (d1) -| node[above left, font=\small] {не} (a3);
  \draw[line] (d1) -| node[above right, font=\small] {да} (a4);
  \draw[line] (a3) -- (a7);
  \draw[line] (a4) -- (a5);
  \draw[line] (a5) -- (a6);
  \draw[line] (a7) -- (a8);
  \draw[line] (a6) |- (a8);
  \draw[line] (a8) -- (d2);
  \draw[line] (d2) -| node[above left, font=\small] {да} (a10);
  \draw[line] (d2) -| node[above right, font=\small] {не} (a9);
  \draw[line] (a10) |- (e);
  \draw[line] (a9) |- (e);
\end{tikzpicture}
\end{adjustbox}
\caption{Дейности при преминаване през пълен цикъл. Изходният текст на модела е в \code{examples/roundtrip.activity} и се превежда с \code{blueprint activity}}
\label{fig:activity}
\end{figure}
